I. Sieci Mobilne i Standardy 802 (Mobile Networks)Charakterystyka i Sprzęt:Najpopularniejszymi typami sieci bezprzewodowych są sieci komórkowe (cellular) oraz sieci sensorowe.  Pasma ISM (433 MHz, 868 MHz, 2.4 GHz) są wykorzystywane do wszelkich rozwiązań zdalnych.  Technologia Software Defined Radio (SDR) pozwala na zaimplementowanie prawie każdego rozwiązania z zakresu radiokomunikacji przy użyciu jednego urządzenia.  Bezpieczeństwo sieci komórkowych (ISM) zależy od mechanizmów protokołowych i sprzętowych/programowych, a nie od samej częstotliwości.  Klasyfikacja Protokołów 802:Standard 802.11 definiuje sieci WLAN (a/b/g/n/ac/ax/be).  Standard 802.15 definiuje sieci WPAN, w tym technologię Bluetooth (802.15.1).  Standard 802.16 to sieć WiMax, dzieląca się na wersję stacjonarną (802.16d) oraz mobilną (802.16e).  II. WiMax (802.16) - Architektura i BezpieczeństwoArchitektura Sieci: Sieć WiMax składa się z co najmniej jednej stacji bazowej (BS) oraz wielu stacji klienckich (SS - Subscriber Station).  Protokoły Zarządzania Kluczami (PKM):Wersja 802.16d wykorzystuje PKMv1 (oparty na RSA), podczas gdy 802.16e używa PKMv2.  Dla autoryzacji stacje posiadają certyfikat X.509, z którego klucz publiczny RSA jest wbudowany w sprzęt.  PKMv1 korzysta z algorytmów DES i 3DES, brakuje w nim uwierzytelniania stacji bazowej (BS) oraz chroni hashe przy pomocy algorytmu SHA-1 (podatnego na kolizje).  PKMv2 poprawia te błędy, stosując tryb CCM (Counter Mode Cipher Block Chaining Message Authentication Code) oraz zastępując HMAC protokołem CMAC.  Klucze i Powiązania Bezpieczeństwa (Security Associations - SA):Wyróżnia się trzy typy SA: Główne (Primary), Statyczne (Static) oraz Dynamiczne (Dynamic).  Połączenie wymaga użycia klucza autoryzacyjnego (AK) przesyłanego przez BS do stacji klienckiej.  Ostateczne szyfrowanie odbywa się za pomocą kluczy TEK (Transmission Encryption Key), których uzyskanie wymaga poprawnej weryfikacji przez HMAC.  III. Standard 802.11 (Wi-Fi) - Tryby, Autoryzacja i AtakiTryby Pracy:Ad hoc: Wykorzystuje protokół DCF (Distributed Coordination Function), będący wariantem CSMA/CA.  Infrastruktura: Sieć korzysta z protokołu PCF (Point Coordination Function), który jest zmechanizowany centralnie na zasadzie odpytywania stacji.  Protokoły Bezpieczeństwa Wi-Fi:WEP: Używa szyfru strumieniowego RC4, wektora inicjującego IV o długości 24 bitów oraz sumy kontrolnej CRC32. Z uwagi na małą przestrzeń IV wektor ten często ulega kolizji.  WPA: Korzysta z szyfru RC4 (klucze 128-bitowe) wraz z protokołem TKIP (Temporal Key Integrity Protocol), dynamicznie zmieniającym klucze.  WPA2/802.11i: Wymaga stosowania CCMP (Counter Mode with Cipher Block Authentication Code Protocol) oraz algorytmu AES.  WPA3: Wymaga min. trybu Counter mode CBC AES (128 bitów) lub GCB AES (192 bity). Częścią tego standardu jest także OWE (Opportunistic Wireless Encryption) do szyfrowania sesji bez uwierzytelniania.  Uwierzytelnianie (802.1x):Standard obejmuje trzy strony: Suplikanta (Klienta), Autentykatora (Punkt dostępowy) i Serwer Uwierzytelniający (np. RADIUS).  Wykorzystuje protokoły EAP, m.in. EAP-TLS (w tunelu SSL/TLS za pomocą certyfikatów), PEAP (najpierw autentykacja serwera, potem klienta), EAP-TTLS, a także EAP-MD5, LEAP, EAP-FAST (Cisco) czy EAP-SIM / EAP-AKA dla sieci GSM/3G.  Wektory Ataków na Wi-Fi:Kradzież sesji i spoofing MAC: Atakujący modyfikuje adres MAC, zmuszając klienta do rozłączenia przy użyciu fałszywych pakietów DEAUTH i wykonując w ten sposób atak Man-in-the-Middle (MitM) z użyciem fałszywego AP.  Łamanie WEP/WPA: Ataki statystyczne i słownikowe (brute force dla krótkich haseł WPA-PSK), jak również ataki bazujące na przewidywalnych wektorach IV.  Ataki DoS: Zalewanie pamięci AP fałszywymi żądaniami autoryzacji (Associate flood).  Ataki na "Sleep Mode": Wysłanie fałszywych żądań wybudzenia by przechwycić pakiety przetrzymywane dla stacji mobilnej.  IV. Transakcje Elektroniczne (Electronic Transactions)Narzędzia i Metody:Transakcje przeprowadzane są za pomocą infrastruktury kluczy publicznych (PKI), gdzie klasyczne szyfrowanie symetryczne (np. AES) zapewnia poufność danych, a algorytmy asymetryczne (np. RSA) wykorzystywane są do uwierzytelniania i podpisów cyfrowych.  Protokoły Transakcyjne:Standard SET (Secure Electronic Transactions): Wprowadza koncept "podwójnego podpisu", który pozwala ukryć szczegóły zamówienia przed bankiem oraz detale karty płatniczej przed sprzedawcą.  3D-Secure: Wymaga dodatkowego potwierdzenia płatności przy użyciu dodatkowego kanału (np. aplikacji mobilnej) dla kart płatniczych.  OTP (Internet Open Trading Protocol): Integruje klasyczne zwyczaje handlowe (negocjacje, formy dostawy) w świecie wirtualnym.  Pieniądze Cyfrowe i Systemy:Wirtualna gotówka powinna zapewniać anonimowość, niezależność, podział, możliwość transakcji offline i ochronę przed podwójnym wydaniem.  ECASH: Monety są chowane i podpisywane "w ciemno" przez bank, by zachować anonimowość zakupów.  Opayo / Sage Pay: Proces zakłada użycie Serwerów Gotówkowych (Currency Servers), które weryfikują certyfikaty w bazie danych.  CAFE / EPI: Rozwiązanie offline oparte o karty inteligentne (smart cards) wyposażone w koprocesor „Guardian” do ochrony przed duplikacją płatności.  Kryptowaluty (Cryptocurrency):Funkcjonują w architekturze P2P i służą jako zdecentralizowana baza danych chroniona protokołami kryptograficznymi (np. Bitcoin, Ethereum).  Adres Bitcoin generowany jest w oparciu o algorytm ECDSA na krzywej secp256k1. 256-bitowy klucz publiczny zostaje przekształcony przez hashe SHA-256 oraz RIPEMD-160, a finalnie zakodowany przy pomocy Base58Check.  Standardy Branżowe:PCI DSS (Payment Card Industry Data Security Standard): Definiuje 12 wymagań weryfikacyjnych dla operatorów płatności w celu ochrony danych kart, w tym: zakaz używania domyślnych haseł, instalację zapór sieciowych (firewall) i regularne testowanie.  


I. Transakcje Internetowe (Electronic Transactions)Definicja: Transakcja to procedura realizowana w formie protokołu przez co najmniej dwie strony, z możliwym udziałem arbitra lub zaufanej strony trzeciej (TTP).  Rodzaje transakcji: Transakcje bazodanowe, wymiana dokumentów, podpisywanie umów, transakcje finansowe (płatności, transfery) oraz e-commerce.  Zarządzanie transakcjami: Obejmuje kontrolę poprawności, rejestrację, sprawdzanie wiarygodności stron (identyfikacja i autentykacja), tworzenie certyfikatów i kluczy oraz nadzór nad procesem. Działania te realizują strony transakcji oraz zaufane strony trzecie.  Narzędzia kryptograficzne:Algorytmy symetryczne (np. AES) zapewniające poufność.  Algorytmy asymetryczne (np. RSA) do podpisów cyfrowych i uwierzytelniania.  Funkcje skrótu (hash) i generatory pseudolosowe.  II. Protokoły PłatnościStandard SET (Secure Electronic Transactions):Stworzony w 1997 r. przez m.in. VISA i MasterCard, oparty na infrastrukturze PKI.  Wykorzystuje podwójny podpis cyfrowy ($S_A(H(H(M1), H(M2)))$), który ukrywa dane karty przed sprzedawcą oraz specyfikację zakupów przed bankiem.  Obecnie wyparty m.in. przez promowany przez VISA standard 3D-Secure.  3D-Secure:Dodatkowy protokół autoryzacji wykorzystujący odrębny kanał (np. aplikację mobilną) do potwierdzania transakcji.  Opiera się na wzajemnym uwierzytelnianiu, komunikatach XML i tunelach SSL.  OTP (Internet Open Trading Protocol):Protokół integracyjny przenoszący klasyczne zwyczaje handlowe do sieci (oferty, negocjacje, formy płatności, paragony).  III. Pieniądze Cyfrowe (Digital Cash)Cechy idealnego systemu: Niezależność od lokalizacji, bezpieczeństwo (ochrona przed kopiowaniem), prywatność (anonimowość), możliwość działania off-line, transferowalność, podzielność, a także skalowalność i akceptowalność.  Historyczne/Przykładowe systemy:ECASH (1994-1998): W pełni anonimowy system, gdzie monety były ukrywane, podpisywane w ciemno przez bank i przesyłane do portfela użytkownika.  Opayo (Sage Pay): Opiera się na serwerach walutowych (Currency Servers), które kontrolują cyfrowe pieniądze pod kątem podwójnego wydania. Pieniądz to podpisany kluczem publicznym komunikat.  CAFE / EPI: Projekt europejski oparty na płatnościach off-line za pomocą kart inteligentnych (smart cards) wyposażonych w koprocesor zabezpieczający "Guardian".  IV. Kryptowaluty (Cryptocurrency)Charakterystyka: Rozproszony system księgowy (P2P), gdzie stan posiadania kontrolowany jest przez właściciela klucza prywatnego.  Generowanie adresu Bitcoin:Wykorzystanie algorytmu ECDSA (standard secp256k1) do wygenerowania 256-bitowego klucza prywatnego i publicznego.  Przekształcenie klucza publicznego przez funkcję SHA-256 oraz RIPEMD-160 (uzyskanie 160-bitowego ciągu).  Dodanie 4 bajtów sumy kontrolnej i bajtu wersji (200-bitowy adres).  Zakodowanie całości algorytmem Base58Check.  V. Bezpieczeństwo Danych Kart (PCI DSS)Standard PCI DSS: Zawiera 12 rygorystycznych wymagań dla operatorów, w tym m.in.: poprawną konfigurację firewalla, zakaz używania domyślnych haseł, szyfrowanie transmisji, ochronę antywirusową, śledzenie zasobów sieciowych i ograniczanie dostępu fizycznego.  VI. Zarządzanie Kluczami i Infrastruktura PKI (Public Key Infrastructure)Procedury zarządzania: Generowanie kluczy, dystrybucja certyfikatów, negocjacja sesji (Diffie-Hellman, protokół STS), przechowywanie (escrowing - przetrzymywanie kluczy przez TTP) oraz unieważnianie.  Modele PKI:X.509: Tworzy globalną strukturę hierarchiczną (drzewo) z korzeniem ROOT CA. Posiada listy unieważnionych certyfikatów (CRL) utrzymywane przez CA. Weryfikacja następuje przez tzw. "łańcuch certyfikatów". Opiera się na systemie nazw X.500.  PGP: Pakiet oprogramowania tworzący sieć "peer-to-peer" (Web of Trust) bez centralnego CA. Użytkownicy lokalnie przechowują i potwierdzają klucze innych osób, jednak problematyczne jest ich unieważnianie.  SPKI / SDSI: Infrastruktury rozproszone oparte na lokalnych przestrzeniach nazw. SPKI skupia się na kontroli dostępu za pomocą certyfikatów autoryzacji (AC) i certyfikatów nazw (NC). Posiada czasowe listy CRL i bazuje na relacjach wyrażonych jako krotki (5-krotka w SPKI). Rolę CA pełnią sami użytkownicy.  

# Bezpieczeństwo systemów komputerowych: Bezpieczeństwo poczty elektronicznej (2026)

**Autor:** PhD E.Eng Piotr Rajchowski

---

## 1. Popularne zagrożenia sieciowe (Internetowe)
Prezentacja opisuje szereg zagrożeń, z jakimi borykają się sieci komputerowe:

* [cite_start]**TCP Spoofing / Sequence Number Guessing:** Przewidywanie sekwencji numerów TCP w celu podszycia się pod zaufany host[cite: 1, 2, 3].
* [cite_start]**Sniffing:** Odczytywanie danych przesyłanych w sieci, dla których nie były one przeznaczone (aktywne i pasywne)[cite: 4].
* [cite_start]**Impersonating (Spoofing):** Podszywanie się pod inne urządzenia w każdej warstwie protokołu IP[cite: 5].
* **Rodzaje spoofingu:**
    * [cite_start]MAC addresses[cite: 6].
    * [cite_start]ARP (Address Resolution Protocol)[cite: 6].
    * [cite_start]DNS (Domain Names System) – przejęcie kontroli nad serwerem DNS[cite: 6].
    * [cite_start]IP routing – wymuszanie błędnej ścieżki pakietów[cite: 7].
    * [cite_start]WWW – fałszywe strony internetowe, nadużycia serwerów proxy[cite: 7].
* [cite_start]**DoS (Denial of Service):** Przeciążanie zasobów systemowych lub sieciowych (np. ataki na WWW, SMTP, DNS)[cite: 9, 10].
* [cite_start]**Port scanning:** Przeszukiwanie portów w poszukiwaniu podatności[cite: 11].
* [cite_start]**Zagrożenia przeglądarek:** Kod HTML, wtyczki i moduły z niepewnych źródeł[cite: 11].

---

## 2. Bezpieczeństwo serwerów WWW i Java
* [cite_start]**Zabezpieczenia Java:** Wykorzystanie kryptografii (RSA, AES, ChaCha20, SHA-3), PKI (X.509, CRLs) oraz JSSE dla SSL/TLS[cite: 12, 13].
* **Bezpieczna konfiguracja serwera WWW:**
    * [cite_start]Usunięcie niepotrzebnych kompilatorów (np. gcc/g++)[cite: 14].
    * [cite_start]Właściwa konfiguracja usług (ostrożność przy FTP, SSH)[cite: 15, 16].
    * [cite_start]Implementacja zabezpieczeń na poziomie aplikacji, serwera (np. dyrektywa Apache `<Directory>`) oraz systemu[cite: 17, 18, 19].
    * [cite_start]Stosowanie firewalla i środowisk chroot[cite: 19].

---

## 3. VPN (Virtual Private Network)
* [cite_start]Służy do tworzenia bezpiecznych kanałów komunikacji w środowiskach niepewnych[cite: 46].
* [cite_start]Technologie: SSL VPN (przeglądarka, OpenVPN) oraz IPSec[cite: 47].

---

## 4. Protokoły SSL/TLS
* [cite_start]**Zadania:** Szyfrowanie, zapewnienie integralności danych i uwierzytelnianie[cite: 48, 49].
* [cite_start]**Struktura:** Złożony z podprotokołów (Handshake, Record, Alert, Cipher Change)[cite: 50].
* [cite_start]**TLS 1.3:** Uproszczony proces uzgadniania (handshake), szybsze nawiązywanie połączenia (0-RTT), wyższe bezpieczeństwo[cite: 70, 81, 83].
* [cite_start]**DTLS:** Wersja dla protokołów datagramowych (UDP)[cite: 87, 88].

---

## 5. Bezpieczeństwo poczty elektronicznej
* [cite_start]**Podstawowe komponenty:** MUA (klient poczty), MSA, MTA (serwer poczty), MDA, MAA[cite: 96, 97, 98, 99].
* [cite_start]**Wymagania:** Poufność, autentyczność, integralność, niezaprzeczalność[cite: 113].
* **Technologie zabezpieczające:**
    * [cite_start]**PGP/GnuPG:** Szyfrowanie i podpisywanie (klucze publiczne/prywatne, pierścienie kluczy)[cite: 114, 121, 156].
    * [cite_start]**PEM (Privacy Enhanced Mail):** Oparte na hierarchii certyfikatów CA i standardzie X.509[cite: 132, 136].
    * [cite_start]**S/MIME:** Następca PEM, wspierający uwierzytelnianie, integralność i poufność (często zintegrowany z klientami poczty jak Outlook)[cite: 158, 159, 164].
    * [cite_start]**DKIM (DomainKeys Identified Mail):** Pozwala organizacji potwierdzić odpowiedzialność za wysłaną wiadomość (podpis w nagłówku, publiczny klucz w DNS), co pomaga w walce z phishingiem i spoofingiem[cite: 174, 175, 178, 179].


# Data Centers Interconnections - Podsumowanie

## 1. Wstęp i Koncepcja
* [cite_start]**Współczesne centra danych:** Zakłada się, że są one rozproszone, a nie scentralizowane[cite: 117].
* [cite_start]**Główny cel:** Zapewnienie niezawodnej i szybkiej wymiany danych między centrami [cite: 118] [cite_start]oraz dostarczanie usług klientom o zróżnicowanych potrzebach przy zachowaniu jakości usług niezależnie od obciążenia[cite: 119].
* **Koncepcja rozproszonych centrów danych:**
    * [cite_start]Wymaga połączenia co najmniej dwóch jednostek[cite: 134].
    * [cite_start]Jednostki aktywnie i dynamicznie równoważą obciążenie[cite: 135].
    * [cite_start]Zapewniają usługi w tle, takie jak adaptacyjna lokalizacja danych czy kopie zapasowe[cite: 136].
    * [cite_start]Cel: minimalizacja opóźnień i maksymalizacja stabilności[cite: 137].

## 2. Klasyfikacja Połączeń
Połączenia klasyfikuje się na trzy grupy:
* [cite_start]**Duże:** np. połączenia podmorskie (międzykontynentalne)[cite: 143].
* [cite_start]**Średnie:** do 100 km, np. połączenia między centrami danych w regionie miejskim[cite: 143].
* [cite_start]**Małe:** w obrębie małych regionów (kilka kilometrów) lub między oddziałami jednego centrum[cite: 144].
* [cite_start]**Typowy standard:** Stosuje się co najmniej dwie niezależne linie (od różnych dostawców lub niezależne)[cite: 154].

## 3. Aspekty Techniczne
* **Physical Layer:**
    * [cite_start]OTN (Optical Transport Network) – pakowanie danych w jeden lub wiele strumieni (do 400 Gb/s)[cite: 169, 170].
    * [cite_start]"Extended" Ethernet (Data center bridging) – minimalizacja utraty pakietów[cite: 171, 172].
    * [cite_start]CWDM (Coarse Wavelength Division Multiplexing) – zwiększanie przepustowości[cite: 173, 174].
* [cite_start]**Architektura sieci:** Ewolucja w kierunku **leaf-spine** w celu poprawy skalowalności i obniżenia opóźnień[cite: 183, 185].
* [cite_start]**ROADM:** Wykorzystywane do połączeń lokalnych lub dzielenia strumieni danych[cite: 213].
* [cite_start]**Benchmarking:** Testowanie przepustowości, opóźnień, jittera i utraty ramek (np. standard RFC 2544)[cite: 214, 215].

## 4. Synchronizacja Sieciowa
[cite_start]Synchronizacja czasu jest kluczowa dla orkiestracji zdarzeń i korelacji danych[cite: 218].
* **Metody synchronizacji:**
    * [cite_start]GNSS (globalna dystrybucja czasu)[cite: 219].
    * [cite_start]Protokoły lokalne: NTP (Network Time Protocol) oraz PTP (Precision Time Protocol)[cite: 220].
* **NTP:**
    * [cite_start]Oparty na hierarchii drzewa[cite: 222].
    * Dokładność ok. [cite_start]100 ms[cite: 230].
    * [cite_start]Zagrożenia: ataki man-in-the-middle, opóźnienia pakietów UDP[cite: 225, 226].
* **PTP (IEEE 1588):**
    * [cite_start]Oparty na architekturze master-slave[cite: 232].
    * [cite_start]Zapewnia dokładność nanosekundową[cite: 231].
    * [cite_start]Zagrożenia: DoS, spoofing, modyfikacja pakietów[cite: 238].

## 5. Zagrożenia GNSS (Jamming/Spoofing)
* [cite_start]**Jamming (zagłuszanie):** Zakłócanie sygnału GNSS (np. za pomocą SDR) w celu uniemożliwienia estymacji pozycji/czasu[cite: 63, 64].
* [cite_start]**Spoofing:** Emisja fałszywego sygnału imitującego konstelację satelitów w celu wywołania błędnej estymacji pozycji/czasu[cite: 65, 70].
* **Obrona:**
    * [cite_start]Galileo OS-NMA (autentykacja wiadomości nawigacyjnych)[cite: 89].
    * [cite_start]Wykorzystanie protokołów takich jak TESLA (lightweight, dla komunikacji stratnej)[cite: 93, 94].
    * [cite_start]Analiza spójności komunikatów (np. NMEA)[cite: 81, 95].


mindmap
  root((Practical Security Issues 2026))
    Slide 1 - Introduction
      Security of Computer Systems 2026
      Practical Security Issues v1.0
      Lecturer Piotr Rajchowski
      Date 10.05.2026

    Slide 2 - AI Agent Environments
      Development of AI agents
      Need for isolated or dedicated data exchange environments
      Example Moltbook
      Analysis based on concept not content
      Framework where
        Entries created by AI agents
        Entries accessed by AI agents
        Entries reused by AI agents

    Slide 3 - OpenClaw Basics
      Moltbook based on OpenClaw
      OpenClaw
        Agent perspective
        Server software perspective
      Not AI model
      Not bot
      WebSocket communication
      Data exchange between services

    Slide 4 - Service Access Model
      Access methods important
      User accesses service through OpenClaw
      OpenClaw accesses LM model
      Skills
        ML code extensions
        Link models with OpenClaw
      Theory
        Local OpenClaw
        External ML models
      Practice
        Both communications often external

    Slide 5 - AI Agent Attacks
      Skills and models usually not attacked directly
      Attack vector
        Prompt injected messages
      Malicious content published
      Message may
        Remain unnoticed
        Gain importance
        Be accessed by other agents

    Slide 6 - Prompt Injection Effects
      Public content contains injected prompts
      AI agent behavior influenced
      Prompts may be plain text
      Malicious agents can manipulate visibility
      Detection of malicious prompts difficult

    Slide 7 - Hidden Malicious Prompts
      Plain text appearance
      No obvious malicious indicators
      Actions appear normal
      Security consequences
        Data exchange manipulation
        Decision biasing
      User-granted privileges increase risk

    Slide 8 - Moltbook Security Issues
      Database access issues
      Authentication keys exposed
      Agents had RW permissions
      Security companies scan
        Skills
        Messages
      Verified OpenClaw libraries planned
      Main challenge
        Threat exists in content
        Not in code

    Slide 9 - Attack Scenario
      Typical user permissions abused
      Hidden prompts inside normal content
      User launches standard task
      Examples
        Code checking
        Code review
        Code submission
        Transactions

    Slide 10 - Agent Compromise
      Agent unaware of malicious prompt
      API communication occurs normally
      Hidden activities executed
      Possible outcomes
        Data exfiltration
        Secondary communication with attacker
        Modified user interactions

    Slide 11 - Major AI Agent Threats
      Prompt injection
      Decision biasing
      Excessive privileges
      Environment threats
      Framework threats
      Input data modification
      Code execution

    Slide 12 - AI Interaction Control
      New control methods required
      Cloudflare AI Labyrinth
      Endless link structures
      Monitoring trigger
      Scraping control
      User decides
        Whether scraping allowed
        How much scraping allowed

    Slide 13 - JA3 Fingerprinting
      Monitor agent and bot access
      JA3 fingerprinting
      Identifies
        Users
        Bots
        Agents
      Based on TLS session characteristics
      Known fingerprints for untrusted entities

    Slide 14 - JA3 Generation
      Based on Client Hello exchange
      Uses
        TLS version
        Elliptic curves
        Cipher identifiers
        Extensions
        Signature algorithms
        Hash algorithms
      Google GREASE
        Random extensions
        Multiple fingerprints
      Compensation mechanisms exist

    Slide 15 - JA4 Introduction
      JA4 newer than JA3
      Different methodology
      Covered later
      Other access-control mechanisms exist
      User action verification possible

    Slide 16 - JA4+ Network Fingerprinting
      JA4
        TLS client fingerprinting
      JA4H
        HTTP client fingerprinting
      JA4X
        X509 certificate fingerprinting
      JA4TScan
        TCP scanner fingerprinting
      JA4D
        DHCP fingerprinting

    Slide 17 - JA4 Practical Use
      Fingerprints for
        Approved systems
        Regular users
        Aggressive bots
      Example Windows 11 JA4D fingerprint
      Traffic discovery plugins available
      Monitoring capabilities

    Slide 18 - Amazon CloudFront
      Proxy server network
      Content Delivery Network
      VPC integration
      Application content routed through VPC
      Reduced public exposure
      Access controlled through CloudFront

    Slide 19 - Cloud Tool Vulnerabilities
      Threats in cloud-connected tools
      Example Copy-Fail
      Linux kernel exploit
      Local privilege escalation
      Root access acquisition

    Slide 20 - Transition Slide
      Introduction to vulnerability details

    Slide 21 - Copy-Fail Technical Details
      Vulnerable component
        algif_aead interface
      Cryptographic module
      AEAD
        Authenticated Encryption with Associated Data
      Input-output handling error
      Page cache access possible
      Privilege escalation to administrator

    Slide 22 - Mitigation Steps
      Limit local access
      Limit physical access
      Restrict SSH access
      Use SELinux
      Use SCC mode
      Apply kernel updates
      Revert vulnerable optimization

    Slide 23 - SELinux Overview
      Security-Enhanced Linux
      Kernel modifications
      Security policy enforcement
      Access control
      Socket restrictions
      File system controls

    Slide 24 - SELinux Modes
      Enforcing
        Block policy violations
      Disabled
        No protection
      Permissive
        Allow violations
        Log all actions

    Slide 25 - Security Context Constraints
      Container platforms
      Pod permission control
      Restrict dangerous resources
      Protect
        File system
        Administrative functions
      Controlled privileged access

    Slide 26 - Kernel Security Updates
      CVE-2026-23136
        libceph DoS
      CVE-2026-23270
        act_ct use-after-free
        DoS or privilege escalation
      CVE-2026-31402
        NFSv4 heap overflow
      CVE-2026-31431
        algif_aead vulnerability

    Slide 27 - Large Service Security
      Continuous monitoring
      Regular kernel compilation
      Lifecycle management
      Risk analysis
      Threat analysis
      Organized response actions

    Slide 28 - Cloudflare Response Part 1
      Map blast radius
      Assess vulnerable versions
      Determine exposure
      Validate detections
      Confirm exploit visibility

    Slide 29 - Cloudflare Response Part 2
      Proactive threat hunting
      Review previous logs
      Runtime mitigation development
      Safe deployment planning
      Kernel rollout preparation

    Slide 30 - Cloudflare Response Part 3
      Continue software updates
      Kernel deployment
      Controlled server reboot process
      Fleet-wide rollout

    Slide 31 - Final Conclusion
      Threat mitigation through kernel updates
      Edge data center protection
      Security maintenance required


